\documentclass[12pt]{article}

% Set Times New Roman equivalent
\usepackage{mathptmx}  % Times font (serif)
\usepackage[T1]{fontenc}
\usepackage[utf8]{inputenc}
\usepackage[english]{babel}

% Page layout settings
\usepackage[margin=1.87cm]{geometry}

% Line spacing
\usepackage{setspace}
\setstretch{1.0}  % single-spaced

% Paragraph formatting
\usepackage{parskip}  % Adds space between paragraphs instead of indent
\usepackage{ragged2e}
\RaggedRight  % Left-justified, not fully justified
\usepackage{float}
\usepackage{graphicx}
\usepackage{fancyhdr}
% For references
\usepackage[square,numbers]{natbib}
\bibliographystyle{plainnat}  % or change to any appropriate style


\pagestyle{fancy}
\fancyhf{} % clear header/footer
\fancyhead[R]{Xiaomao Wang} % name on right side
\fancyfoot[C]{\thepage}
\renewcommand{\headrulewidth}{0pt}

\begin{document}
\section*{Part I – Contributions to research and development}
\textbf{a. Articles published or accepted in peer-reviewed journals}\\
\medskip
Tomlinson, K. W., Yu, F., \textbf{Wang, X.,} Yao, X., Yu, C. C., Charles‐Dominique, T., ... \& Gélin, U. (2025). The macroecology of spines on woody plants. \textit{Biological Reviews}. \\
\medskip
\textbf{Wang, X.,} Fan, H., Phoncharoen, W., Gélin, U., \& Tomlinson, K. W. (2023). Leaf chemistry of architecturally defended plants responds more strongly to soil phosphorus variation than non‐architecturally defended ones. \textit{Physiologia Plantarum, 175}(1), e13856.\\
\medskip
\textbf{Wang, X.,} \& Ibáñez, I. (2022). The contrasting effects of local environmental conditions on tree growth between populations at different latitudes. \textit{Forests, 13}(3), 429 (Master work)\\
\medskip
\textbf{b. Other peer-reviewed contributions}\\
\textbf{c. Non-peer-reviewed contributions}\\
\medskip
Wang X.* (2022) Leaf chemistry of architecturally defended plants responds more strongly to soil phosphorus variation than non-architecturally defended ones, oral presentation at the 2022 Xishuangbanna Tropical Botanical Garden Annual Academic Conference held in Xishuangbanna Tropical Botanical Garden, Chinese Academy of Science, Xishuangbanna, China\\
\medskip
Wang X.*, Wang H.,\& Liao A. (2018) Environmental Characteristics of Seagrass Beds \textit{Halophila beccarii} at
Tangjiawan, Zhuhai, poster presentation at the Sixth Undergraduate Science Technology Poster Presentation Conference at BNU-HKBU United International College, Zhuhai, China\\
\medskip
Wang X.*, Roybal O., Hospeti L.,\& Lenihan M. (2017) The influence of sterilization on germination of six \textit{Asclepia} species, poster presentation at the 11th Annual Summer Undergraduate Research Conference at Trinity University, San Antonio, USA\\
\medskip
\textbf{d. Technology transfer}\\
\textbf{e. Contributions resulting from your participation in industrially relevant R\&D activities}\\
\textbf{f. Awarded and submitted patents and copyrights}\\
\section*{Part II – Most significant contributions to research and development}
\textbf{Investigate climate impact on growth} My master’s project focused on how environmental changes influence the growth of the same species across different latitudes. Our results showed that even populations located near the northern edge of their distribution are already experiencing climate-related stress and may struggle to persist. We used models to predict future performance under climate change, and the sugar maples are projected to decline. This work, published in \textit{Forests}, emphasized the importance of spatial higher-resolution in predicting species responses to climate change. This study also helped me transition smoothly into my PhD research by helping me build fundamental skills in Bayesian approach and dendrochronology.\\
\textbf{Collaborative work in macroecology} As a research assistant to Dr. Kyle Tomlinson, I contributed to a global study on the distribution and evolution of spiny plants. I led data cleaning and spine identification and engaged with international collaborators, including contacting experts to verify spine morphologies. This experience strengthened my ability to manage complex datasets, collaborate in a big scale, and contribute meaningfully to spiny plants research.\\
\textbf{Promote endangered seagrass conservation} My undergraduate senior project focused on studying the only known patch of endangered \textit{Halophila beccarii} seagrass in Zhuhai, which at the time was not under any protection. Towards the end of my project, the local government proposed a beach restoration plan that involved relocating sand to the original rocky shoreline, which directly threatens the seagrass habitat. We sent the central government a letter highlighted the ecological importance of the area, supported with field data showing higher biodiversity in seagrass bed compared to nearby bare ground. After collaborating with research teams organized by the government, the seagrass patch has now been officially designated as a protected area with my help. This made me passionate about conservation and helped me learn how to communicate science to policy makers.\\
\section*{Part III – Applicant’s statement}
\textbf{Research Experience}\\
I started my ecology journey studying plant communities and their environmental interactions during my undergraduate senior project, where I investigated the community structure of an endangered seagrass, \textit{Halophila beccarii}. This work revealed how seagrass competes with algae under different  environmental and nutrient conditions. Through this project, I gained hands-on experience in fieldwork planning, data collection, and basic statistical analysis. These experiences sparked my deeper interest in plant-environment relationships.\\
Expanding my focus from intertidal zones to terrestrial ecosystems, my master’s research explored forest dynamics under climate change. My thesis examined how temperature and precipitation influence the annual growth of two common maple species at different latitudes in Michigan, using dendrochronology and Bayesian modeling. This work provided valuable training in tree-ring analysis, advanced statistical modeling in R, and ecological interpretation of long-term climate data.\\
Following my master’s degree, I worked as a research assistant in a community ecology lab. My responsibilities included assisting students with their fieldwork, sorting and identifying plant specimens, data cleaning and analysis. This role further strengthened my skills in R, data management, and collaborative research.\\
My cumulative experiences spanning diverse ecosystems, field and lab techniques, and quantitative tools, confirmed my passion for research and prepared me for a qualified PhD. I started my doctoral program with a strong foundation in ecological research and statistical knowledge, which made me ready to further explore climate change impacts on forest dynamics.\\
\textbf{Relevant activities}\\
I consider myself a well-trained ecologist with a strong foundation in both fieldwork and theory. During my second year in undergrad, I was the only junior student selected for the opportunity to go to Texas for a fully funded summer research program. During my master’s, I took additional coursework to deepen my knowledge across various ecological disciplines and gained extensive hands-on experience through fieldwork in different ecosystems. I also developed strong statistical skills by participating in multiple workshops on advanced statistics and Bayesian modeling, and applying these skills to different ecological datasets.  In recognition of this academic and research background, I was awarded the Faculty of Forestry Graduate Award and the President's Academic Excellence Initiative PhD Award at the start of my PhD.\\
I always try to help younger generation students with research interests. While working as a research assistant, I organized free R workshops to help undergraduate students build coding skills. During my PhD, I also trained and mentored several WorkLearn students, guiding them through both my projects and their own projects. As an international student, I understand the potential inequities in research and always support broader participation and equity whenever I have a chance.\\
Beyond academia, I actively engage in outreach and volunteering. I worked with the Zhuhai Bird Watching Association,  to organize public education activities on seagrass habitats. During my PhD, I also volunteered for the EcoEvo Annual Retreat reaching out to speakers and organized the talks.

\end{document}